% -----------------------------------------------------------------------------
% This file is automatically generated.
% Do not edit manually.
% Source: notebooks/support/regression-analysis/support.Rmd
% -----------------------------------------------------------------------------

\par\addvspace{\topsep}
\begingroup
\fvset{listparameters={%
  \setlength{\topsep}{0pt}%
  \setlength{\partopsep}{0pt}%
  \setlength{\parsep}{0pt}%
  \setlength{\itemsep}{0pt}%
}}
\begin{Verbatim}[breaklines=true,formatcom=\color{ada_blue}]
# Initialize a data frame to store the simulated regression lines
simulated_lines <- data.frame()

# Repeat the process `rep` times to simulate multiple regression lines
for (i in 1:rep) {
  # Sample n points from the relationship y = beta_0 + beta_1 * x + epsilon
  epsilon <- rnorm(n, mean = 0, sd = sigma) # Random noise
  y_values <- beta_0 + beta_1 * x_values + epsilon

  # Fit a linear model
  model <- lm(y_values ~ x_values)

  # Predict fitted values for the extended range of x
  predicted_values <- predict(model, newdata = data.frame(x_values = x_extended))

  # Store the fitted line over the extended x range
  simulated_lines <- rbind(
    simulated_lines,
    data.frame(
      x = x_extended,
      y = predicted_values,
      rep = i
    )
  )
}
\end{Verbatim}
\endgroup
\par\addvspace{\topsep}

\noindent All the elements created earlier are now plotted on a single canvas.
